\answer{ex:14:1}{(а)~На $0.17$~с и на $1.5$~с. (б)~С расстояния меньше $11$~см.}
\answer{ex:14:2}{Касание без боли не проходит ни при каком $\mu$, пока $g\le\beta w_{q\mu}/w_{z\mu}=5$; касание $\mu=1$ проходит при $g>6$: сильная сенсибилизация делает болью обычное прикосновение.}
\answer{ex:14:3}{(а)~$g^*=(1-k_sC)/(1-k_sA)$, устойчиво при $k_sA<1$. (б)~Без касания при $\nu\ge2$, с касанием --- уже при $\nu\ge1.7$.}
\answer{ex:14:4}{$V(t)=-Pe^{-(T-t)/\tau_\gamma}$. (а)~$-Pe^{-T/\tau_\gamma}\approx-0.37P$. (б)~$+Pe^{-(T-t_e)/\tau_\gamma}\approx+0.61P$; он растёт с $t_e$ до $+P$ и учит избегать боли.}
\answer{ex:14:5}{При $k_s=4$ защёлкивания нет (второе устойчивое состояние есть при $k_s\ge4.58$); $T_{\min}\approx4.35$, $1.40$, $0.65$, $0.32\,\tau_s$ при $k_s=4.6$, $5$, $6$, $8$.}
\answer{ex:14:6}{$w_{q\nu}=4/9\approx0.444$, $q_0=2/9\approx0.222$, $w_{q\mu}=49/45\approx1.089$; касание без боли безопасно до $g=49/9\approx5.4$.}
\answer{ex:14:7}{(а)~Порог $\nu_c(g)=\theta/g$ при $\nu_c\ge q_0/w_{q\nu}$, иначе $(\theta+\beta q_0)/(g+\beta w_{q\nu})$; $g^*\approx2.45$, $1.0$, $0.25$. (б)~Открытый вопрос.}
